\documentclass[conference]{IEEEtran}
\IEEEoverridecommandlockouts

\usepackage{cite}
\usepackage{amsmath,amssymb,amsfonts}
\usepackage{algorithmic}
\usepackage{graphicx}
\usepackage{textcomp}
\usepackage{xcolor}
\usepackage{booktabs}
\usepackage{array}
\usepackage{url}
\def\BibTeX{{\rm B\kern-.05em{\sc i\kern-.025em b}\kern-.08em
    T\kern-.1667em\lower.7ex\hbox{E}\kern-.125emX}}

\begin{document}

\title{ChatDBG-MTP: Extending AI-Assisted Debugging with Post-Mortem Analysis,
Repository-Aware Context, and Automated Repair}

\author{
\IEEEauthorblockN{Siddhi Tiwari}
\IEEEauthorblockA{
\textit{Department of Computer Science \& Engineering}\\
\textit{[Your Institution]}\\
[City, Country]\\
siddhitiwari0109@gmail.com
}
}

\maketitle

\begin{abstract}
Debugging is one of the most time-consuming activities in software development.
Existing AI-assisted debugging tools, including the original ChatDBG, require an
active debugger session, limiting their use to crashes reproducible in a live
environment. In this work we present \textbf{ChatDBG-MTP}, an extended version of
ChatDBG that introduces five major contributions: (1)~\textbf{post-mortem analysis}
that runs on crash log files without a live process, (2)~a \textbf{zero-friction
exception hook} that auto-triggers analysis on any unhandled exception,
(3)~\textbf{repository-aware context} that follows import graphs to provide the LLM
with a broader view of the codebase, (4)~\textbf{git-aware context} that surfaces
recent commit history, blame information, and diffs to support regression
detection, and (5)~\textbf{automated fix application and test generation} that
complete the full repair cycle. We evaluate ChatDBG-MTP on 8 benchmark programs
across four configuration variants. Our results show that adding repository
context improves root-cause keyword coverage from 57\% to 77\% (+20~pp),
fix correctness from 50\% to 75\% (+25~pp), and test generation from 12\%
to 75\% (+63~pp). Combining repository and git context achieves the best
overall accuracy (89\% keyword coverage, 88\% fix correctness) with a
modest cost increase of 38\%.
\end{abstract}

\begin{IEEEkeywords}
debugging, large language models, post-mortem analysis, automated program repair,
fault localization, git-aware debugging, test generation
\end{IEEEkeywords}

%-----------------------------------------------------------------------
\section{Introduction}
%-----------------------------------------------------------------------

Modern software crashes in many contexts where a live debugger is unavailable:
production environments, continuous integration pipelines, remote machines, and
crash reports submitted by end users. Even when a debugger is available, AI
assistants such as ChatDBG~\cite{chatdbg} only analyze the immediate stack trace,
missing broader codebase context that could accelerate diagnosis.

We identify three key gaps in existing AI-assisted debugging:

\begin{enumerate}
\item \textbf{Availability} --- tools require a live debugger session, excluding
  the majority of real-world crashes captured in logs.
\item \textbf{Context} --- analysis is limited to the call stack, ignoring imported
  modules, class hierarchies, and recent code changes that frequently contain
  the root cause.
\item \textbf{Completeness} --- tools explain bugs but do not close the repair
  loop: applying the fix and writing a regression test.
\end{enumerate}

ChatDBG-MTP addresses all three gaps through the following contributions:
\begin{itemize}
  \item A \textbf{post-mortem analysis} mode (\texttt{chatdbg --analyze crash.log})
    that parses Python tracebacks from text files and runs LLM analysis without
    any running process.
  \item A \textbf{zero-friction exception hook} (\texttt{chatdbg.install()}) that
    registers a \texttt{sys.excepthook} so every unhandled exception is
    automatically analyzed.
  \item A \textbf{repository context} mechanism (\texttt{--repo}) that traverses
    the import graph from crashed frames to supply relevant source files.
  \item A \textbf{git context} mechanism that automatically includes
    \texttt{git log}, \texttt{git blame}, and \texttt{git diff} output for each
    crashed file.
  \item An \textbf{automated repair} pipeline using two LLM tools:
    \texttt{apply\_fix} and \texttt{generate\_test}.
  \item An \textbf{evaluation harness} for systematic comparison of configuration
    variants.
\end{itemize}

%-----------------------------------------------------------------------
\section{Background and Related Work}
%-----------------------------------------------------------------------

\subsection{ChatDBG}

ChatDBG~\cite{chatdbg} integrates an LLM into the Python \texttt{pdb}/\texttt{ipdb}
debugger and the native GDB/LLDB debuggers. When the user types \texttt{why},
ChatDBG sends the stack trace, error message, and local variable values to the
LLM, which can call back into the debugger via registered functions
(\texttt{debug}, \texttt{info}, \texttt{slice}). ChatDBG uses \texttt{litellm}
to support multiple model providers.

\subsection{Automated Program Repair}

APR tools such as GenProg~\cite{genprog}, Prophet~\cite{prophet}, and
Angelix~\cite{angelix} use search or symbolic analysis to generate patches.
Recent LLM-based approaches (e.g., AlphaRepair, ChatRepair) prompt models with
the buggy code and test suite. Unlike these tools, ChatDBG-MTP operates at
debugging time on a crashed program, without requiring a pre-existing test suite
or formal specification.

\subsection{Fault Localization}

Spectrum-based fault localization (e.g., Tarantula~\cite{tarantula},
Ochiai) ranks suspicious statements based on test pass/fail patterns. Our
approach is complementary: rather than ranking statements, we provide the LLM
with rich context and rely on its reasoning to identify the root cause.

\subsection{LLM Code Understanding}

Recent work shows that LLMs can reason about code at multiple levels of
abstraction~\cite{codex}. Providing additional context---function signatures,
docstrings, related code---consistently improves accuracy~\cite{retrieval_prompts}.
Our \texttt{--repo} feature operationalizes this by automatically surfacing the
most relevant context via import graph traversal.

%-----------------------------------------------------------------------
\section{System Design}
%-----------------------------------------------------------------------

\subsection{Architecture Overview}

ChatDBG-MTP extends the original ChatDBG pipeline with a new post-mortem pathway
that runs independently of any live debugger. The pipeline parses the traceback,
builds source/repo/git context, constructs a prompt, and invokes the LLM
assistant, which may call back into \texttt{apply\_fix} and
\texttt{generate\_test}. The interactive pdb mode and post-mortem mode share
the same \texttt{Assistant} class, listeners, and YAML logging infrastructure.

\subsection{Post-Mortem Analysis}

The \texttt{postmortem} package contains three components:

\textbf{Parser} (\texttt{parser.py}): Parses Python tracebacks into a
\texttt{CrashReport} dataclass containing \texttt{FrameInfo} records (filename,
line number, function, source line) and the exception type/message. Handles
chained exceptions, log files with embedded tracebacks, and exceptions
without messages.

\textbf{Context builder} (\texttt{context.py}): For each user-code frame,
reads \texttt{context\_lines} lines around the crash point from disk using
\texttt{linecache}. Library frames (\texttt{site-packages}, stdlib) are
filtered out.

\textbf{Analyzer} (\texttt{analyze.py}): Orchestrates the pipeline, creates
the \texttt{Assistant} with \texttt{apply\_fix} and \texttt{generate\_test}
tools, runs the LLM query, and logs results.

\subsection{Exception Hook}

\texttt{chatdbg.install()} registers a \texttt{sys.excepthook} that:
(1)~lets Python print the original traceback; (2)~formats the traceback as
text using \texttt{traceback.format\_exception}; (3)~calls
\texttt{analyze\_crash\_text()} in memory; (4)~silently skips
\texttt{SystemExit} and \texttt{KeyboardInterrupt}; and (5)~catches all
ChatDBG errors so they never mask the original exception.

\subsection{Repository Context}

\texttt{repo.py} traverses the local import graph starting from the files in
the traceback. It collects all \texttt{.py} files under \texttt{--repo},
seeds with files directly referenced in the crash frames, then parses
\texttt{import} and \texttt{from ... import} statements using the \texttt{ast}
module to find locally-defined dependencies to depth~2. Discovered files are
formatted into the prompt within a token budget of~$\approx6{,}000$ tokens.

\subsection{Git-Aware Context}

\texttt{git\_context.py} runs three \texttt{git} commands per unique user-code
file via \texttt{subprocess}, as summarized in Table~\ref{tab:git_cmds}.
All commands are timeout-protected (10\,s) and fail silently if git is
unavailable.

\begin{table}[htbp]
\caption{Git Commands Used for Each Crashed File}
\label{tab:git_cmds}
\begin{center}
\begin{tabular}{ll}
\hline
\textbf{Command} & \textbf{Purpose} \\
\hline
\texttt{git log --oneline -8 <file>} & Recent commit history \\
\texttt{git blame -L <s>,<e> <file>} & Last author/date of crash lines \\
\texttt{git diff HEAD\textasciitilde1 -- <file>} & Changes in most recent commit \\
\hline
\end{tabular}
\end{center}
\end{table}

\subsection{Automated Repair}

Two LLM tool functions complete the repair cycle:

\textbf{\texttt{apply\_fix(filename, old\_code, new\_code)}}: Shows a unified
diff and prompts the user \texttt{[y/N]} before writing. Old-code matching
uses line-ending-normalized comparison; CRLF is restored on write.

\textbf{\texttt{generate\_test(filename, test\_code)}}: Generates a pytest
test that reproduces the original crash. The test is designed to fail against
the buggy code and pass after the fix, creating a regression guard.

Both tools bypass \texttt{input()} prompts in eval mode
(\texttt{CHATDBG\_EVAL=1}), capturing proposed changes in memory for the
evaluator to score.

%-----------------------------------------------------------------------
\section{Evaluation}
%-----------------------------------------------------------------------

\subsection{Experimental Setup}

We evaluate ChatDBG-MTP on 8 benchmark programs covering common Python bug
categories (Table~\ref{tab:benchmarks}). We test four configuration variants:
\textbf{Baseline} (traceback + source context), \textbf{+Repo}
(+ import-graph files), \textbf{+Git} (+ git log/blame/diff), and
\textbf{+Repo+Git} (all context). All experiments use GPT-4o via the OpenAI API.

\begin{table}[htbp]
\caption{Benchmark Programs}
\label{tab:benchmarks}
\begin{center}
\begin{tabular}{lll}
\hline
\textbf{ID} & \textbf{Bug Type} & \textbf{Exception} \\
\hline
\texttt{testme\_zero\_division}   & Loop variable starts at 0          & \texttt{ZeroDivisionError}  \\
\texttt{sample\_swapped\_args}    & NumPy arguments swapped            & \texttt{TypeError}          \\
\texttt{off\_by\_one\_index}      & List access beyond bounds          & \texttt{IndexError}         \\
\texttt{none\_attribute\_access}  & Method call on \texttt{None}       & \texttt{AttributeError}     \\
\texttt{type\_mismatch}           & String + integer addition          & \texttt{TypeError}          \\
\texttt{key\_error\_dict}         & Missing dictionary key             & \texttt{KeyError}           \\
\texttt{recursion\_overflow}      & No base case in recursion          & \texttt{RecursionError}     \\
\texttt{value\_out\_of\_range}    & Invalid enum/range value           & \texttt{ValueError}         \\
\hline
\end{tabular}
\end{center}
\end{table}

\subsection{Metrics}

We measure: (1)~\textbf{Exception Match} --- does the response name the expected
exception? (2)~\textbf{Keyword Coverage (KW\%)} --- fraction of expected
root-cause keywords found in the response; (3)~\textbf{Fix Proposed} --- was
\texttt{apply\_fix} called? (4)~\textbf{Test Generated} --- was
\texttt{generate\_test} called? (5)~\textbf{Fix Correct} --- does the proposed
fix make the script exit with code~0? (6)~\textbf{Cost} (USD); and
(7)~\textbf{Tokens} per analysis.

\subsection{Results}

Table~\ref{tab:results} shows aggregate metrics averaged across all 8 benchmarks.

\begin{table}[htbp]
\caption{Aggregate Evaluation Results (8 Benchmarks)}
\label{tab:results}
\begin{center}
\begin{tabular}{lrrrrrrr}
\hline
\textbf{Variant} & \textbf{Exc} & \textbf{KW\%} & \textbf{Fix} & \textbf{Test} & \textbf{FixOK} & \textbf{Cost} & \textbf{Tok} \\
\hline
Baseline     & 88\% & 57\% & 75\%  & 12\%          & 50\%          & \$0.019 & 1641 \\
+Repo        & 100\% & 77\% & 100\% & 75\%          & 75\%          & \$0.024 & 2081 \\
+Git         & 100\% & 72\% & 100\% & 50\%          & 62\%          & \$0.024 & 2069 \\
+Repo+Git    & 100\% & \textbf{89\%} & 100\% & \textbf{88\%} & \textbf{88\%} & \$0.026 & 2511 \\
\hline
\end{tabular}
\end{center}
\end{table}

Key findings:

\textbf{Repository context is the single most impactful feature}, improving
keyword coverage by 20~pp and fix correctness by 25~pp over baseline. This
confirms that root causes frequently lie in imported modules not directly
visible in the stack trace.

\textbf{Git context alone improves recall modestly} (+15~pp KW, +12~pp fix
correctness), with its largest gains on regression-type bugs where the
\texttt{git diff} identifies the introducing commit.

\textbf{Combining both contexts achieves the best results} across all metrics.
The improvement is super-additive for test generation (12\% $\to$ 88\%),
suggesting both context types are needed for the LLM to generate runnable tests.

\textbf{Cost is modest and predictable}: the full +Repo+Git variant costs on
average \$0.026 per analysis, a 38\% increase over baseline for a 56\%
improvement in keyword coverage.

%-----------------------------------------------------------------------
\section{Discussion}
%-----------------------------------------------------------------------

\subsection{When Does +Repo Help Most?}

Repository context provides the largest improvement on bugs where the root cause
is in a helper function imported by the crashed module (e.g.,
\texttt{sample\_swapped\_args}, \texttt{none\_attribute\_access}). For
self-contained bugs, the improvement is marginal. This suggests a heuristic:
use \texttt{--repo} when the crash occurs in application code that imports
project-defined modules.

\subsection{When Does +Git Help Most?}

Git context is most valuable for regression bugs introduced by a recent commit.
In our benchmark, \texttt{off\_by\_one\_index} and \texttt{recursion\_overflow}
showed the largest gains with +Git, consistent with bugs introduced by a code
change visible in \texttt{git diff HEAD\textasciitilde1}. For bugs that have
existed in the codebase for many commits, git context adds noise rather than
signal.

\subsection{Limitations}

\textbf{Static keyword matching} is a proxy for correctness. Future work should
use LLM-as-judge evaluation for semantic correctness.

\textbf{Fix correctness} verifies the script exits cleanly on the original
input but does not guarantee absence of regressions on other inputs.

\textbf{Source file availability}: post-mortem analysis requires source files
to be present at the paths recorded in the traceback.

\textbf{Single-model evaluation}: all experiments use GPT-4o. Future work
should benchmark across Claude, Gemini, and open-source models.

%-----------------------------------------------------------------------
\section{Conclusion}
%-----------------------------------------------------------------------

We presented ChatDBG-MTP, an extension of ChatDBG that enables post-mortem
debugging from crash logs, enriches LLM context with import-graph traversal
and git history, and closes the repair cycle with automated fix application and
test generation. Our evaluation on 8 benchmarks demonstrates consistent,
measurable improvements: adding repository and git context together raises
root-cause keyword coverage from 57\% to 89\% and fix correctness from 50\%
to 88\%, at a modest cost increase of 38\%.

The evaluation harness introduced in this work enables systematic comparison of
context strategies and model choices, providing a foundation for future research
into AI-assisted program repair.

\section*{Acknowledgment}

The author thanks the original ChatDBG team (Berger, Freund, Levin, van Kempen)
for the open-source codebase on which this work builds.

\begin{thebibliography}{00}

\bibitem{chatdbg}
E.~D. Berger, S.~Freund, K.~Levin, and N.~van Kempen,
``ChatDBG: An AI-Powered Debugging Assistant,''
in \textit{Proc. ACM Int. Conf. Foundations of Software Engineering (FSE)},
2025.

\bibitem{genprog}
C.~Le Goues, T.~Nguyen, S.~Forrest, and W.~Weimer,
``GenProg: A Generic Method for Automatic Software Repair,''
\textit{IEEE Trans. Softw. Eng.}, vol.~38, no.~1, pp.~54--72, 2012.

\bibitem{prophet}
F.~Long and M.~Rinard,
``Automatic Patch Generation by Learning Correct Code,''
in \textit{Proc. ACM SIGPLAN-SIGACT Symp. Principles of Programming Languages (POPL)},
pp.~298--312, 2016.

\bibitem{angelix}
S.~Mechtaev, J.~Yi, and A.~Roychoudhury,
``Angelix: Scalable Multiline Program Patch Synthesis via Symbolic Analysis,''
in \textit{Proc. IEEE/ACM Int. Conf. Software Engineering (ICSE)},
pp.~691--701, 2016.

\bibitem{codex}
M.~Chen, J.~Tworek, \textit{et al.},
``Evaluating Large Language Models Trained on Code,''
\textit{arXiv preprint arXiv:2107.03374}, 2021.

\bibitem{retrieval_prompts}
N.~Nashid, M.~Sintaha, and A.~Mesbah,
``Retrieval-Based Prompt Selection for Code-Related Few-Shot Learning,''
in \textit{Proc. IEEE/ACM Int. Conf. Software Engineering (ICSE)},
pp.~2450--2462, 2023.

\bibitem{tarantula}
J.~A. Jones and M.~J. Harrold,
``Empirical Evaluation of the Tarantula Automatic Fault-Localization Technique,''
in \textit{Proc. IEEE/ACM Int. Conf. Automated Software Engineering (ASE)},
pp.~273--282, 2005.

\end{thebibliography}

\end{document}
